\scanpage{129}
% archive sheet number: 128
% top marginal question
\margin{\handwritten{\uline{NB}} \uline{Question} Sous réserve $B$ associative unitaire, si $A$ est commutative, $B$ l'est-elle aussi ? (On suppose $A$ associative unitaire). Question analogue quand on a des structures graduées avec anticommutativité. Donc, la condition de \deleted{\unsure{symé}} comp. aux sym. pour un $c$ est-elle surabondante, en prés. des autres \unsure{compat.|compa\handwritten{t.}} ?}
% source lines 1--3
de base fidèlement plat $k'\handwritten{\longrightarrow}B$, \deleted{\unsure{le foncteur|le torseur}} il devient trivial, donc on a un ``torseur'' sous le monoïde $\unsure{\mathbb{H}|\mathbb{G}}_{k'}$ ; alors le nouveau foncteur fibre $G$ est déduit de $F'=F\otimes_k k'$ par torsion à l'aide de $G$ \ldots

% source lines 4--14
Mais supposons seu\deleted{\unsure{x}}lement $A$ algèbre associative unitaire, et $B$ itou, sans condition de commutativité.\inserted{On suppose $A$ avec \uline{antipodisme}\inserted{\handwritten{$\underline{a}$}} (s'en passer ?).} Donnons-nous un $\handwritten{\sigma\in}U\otimes_k U$ \deleted{\unsure{ci-dessus}} comme ci-dessus pour introduire une relation de commutativité $\handwritten{\psi}$ . Notons (qu'à cela près que $k'$ n'est plus supposé un corps) on trouve une catégorie $C'=C\otimes_k k'$ sur $k'$, avec $\otimes$, définie\inserted{\handwritten{(soit par $U'$, \uuline{soit} \unsure{dual.|aussi.})}} par le ``proanneau'' $V=G(U)=L\overset{\handwritten{\scriptscriptstyle\wedge}}{\otimes}_U(U,\mathrm{int})$, où $(U,\mathrm{int})$ désigne $U$ muni de la structure de $U$-module \deleted{\unsure{(donc de pro-objet de}} (et même d'objet``pro-anneau'' (pour $\otimes$) de $C$) défini à l'aide de sa structure de bimodule sur $U^o\overset{\handwritten{\scriptscriptstyle\wedge}}{\otimes} U$ \deleted{\unsure{à l'aide de}} par restriction des scalaires via $U\handwritten{\xrightarrow{\Delta}}U\overset{\handwritten{\scriptscriptstyle\wedge}}{\otimes} U\handwritten{\xrightarrow{a\otimes\mathrm{id}_U}}U^o\overset{\handwritten{\scriptscriptstyle\wedge}}{\otimes} U$ : Tordant par $L$ et utilisant le fait que l'on a une compatibilité avec $\otimes$, on trouve un pro-anneau $V$ sur $k'$, et ce dernier hérite de $U$ une application diagonale \ldots
\margin{?}
% left margin, a), alongside source lines 9--16
\margin{\handwritten{a) donner un sens intrinsèque à $C\otimes_k k'$, indépendamment de \deleted{$\unsure{\otimes|\circ}$} la donnée de $F$ (\uline{NB} si $k'$ n'est pas un corps, ce n'est pas une catégorie \uline{abélienne}, seulement additive)}}
% source lines 15--29
Ceci\deleted{x} dit, si on a une donnée $\handwritten{\sigma\in}U\overset{\handwritten{\scriptscriptstyle\wedge}}{\otimes} U$ satisfaisant les conditions (1)\inserted{plus haut,} peut-on définir en termes de $\handwritten{\sigma}$ (ou existe-t-il tout au moins) un élément $\handwritten{\tau\in}V\overset{\handwritten{\scriptscriptstyle\wedge}}{\otimes}_k V$ satisfaisant les relations analogues à (1), de telle façon que $(G,c)$ soit compatible aux isomorphismes de symétrie $\handwritten{\psi}$ dans $C$ défini par $\handwritten{\sigma}$, $\handwritten{\psi'}$ dans $C'$ défini par $\handwritten{\tau}$ ??? La seule chance de \deleted{\unsure{montrer}} construire un $\handwritten{\tau}$ canoniquement semble de prouver que $\handwritten{\tau}$ est ``invariant'' lorsqu'on fait opérer $U$ sur $U\overset{\handwritten{\scriptscriptstyle\wedge}}{\otimes}_k U=(U,\mathrm{int})\overset{\handwritten{\scriptscriptstyle\wedge}}{\otimes}(U,\mathrm{int})$ grâce à son \deleted{\unsure{s}}action dite plus haut sur les deux facteurs, ce qui s'exprime par la relation $U^+ .\handwritten{\tau}=0$.\inserted{\handwritten{[Vérifier dans le cas gradué ce qui en est !]}} - Si ceci marchait (ce qui améliorerait le résultat souhaité en NB au haut de la page), cela montrerait que dans l'étude des foncteurs fibres, il y a lieu de laisser tomber les questions de compatibilité aux symétries, celles-ci se récupérant toujours après coup.\inserted{\handwritten{(Mais cf exemples n\textsuperscript{o}12).}}
% left margin, b), alongside source lines 17--21
\margin{\handwritten{b) Montrer que structure $\otimes$, éventuellement avec données $\varphi,\psi,1,\vee$. se prolonge automatiquement à $C'$, avec \unsure{ses axiomes}}}
% left margin, c), alongside source lines 22--30
\margin{\handwritten{c) généraliser, pour un foncteur comme $G:\unsure{C|\mathcal{C}}\to\mathrm{Mod}(k')$, la description du n\textsuperscript{o}6 en description de $C'$ en termes d'un pro-anneau $U'=\mathrm{End}(G)$, et expliciter comment on a encore les diverses structures sur $C$: $\otimes,\varphi,\psi,1,\vee$}}

% source lines 30--31; word continues on scan 130
Il resterait alors la question d'une description \uline{intrinsèque} d'une catégorie $C$ comme dans le n\textsuperscript{o}1,\inserted{\handwritten{avec $\unsure{\check{X}|\check{\otimes}}$}} \unsure{indépendamment|indépendammment} de la donnée d'un $F$ préli%
